\subsection{连续线性映射的余范数}

设 E 和 F 是赋范空间，\(u: E \rightarrow F\) 是连续线性映射，N 是 u 的核，I 是它的像。用 p 表示从 E 到 E/N 上的典范满射，用 i 表示从 I 到 F 中的典范单射。用 \(\widetilde{u}\) 表示从 E/N 到 I 上的双射线性映射，它满足 \(u = i \circ \widetilde{u} \circ p\) 。向量空间 E/N 赋有商范数，即

\[
\| y \| = \inf _ {x \in p^{-1}(y)} \| x \|\tag{2}
\]

对一切 \(y \in E/N\) 成立。映射 \(\widetilde{u}\) 连续且 \(\|u\| = \|\widetilde{u}\|\) ，由此

\[
\| u\| = \sup_{\substack{y\in \mathrm{E / N}\\ y\neq 0}}\frac{\|\widetilde{u} (y)\|}{\|y\|},\tag{3}
\]

其中上确界在 \(\overline{R}_{+}\) 中取。u 的余范数是指

\[
((u)) = \inf_{\substack{y\in \mathrm{E} / \mathrm{N}\\ y\neq 0}}\frac{\| \widetilde{u}(y)\|}{\|y\|},\tag{4}
\]

其中下确界在 \(\overline{R}_{+}\) 中取。因此我们有

\[
((u)) \| y \| \leqslant \| \widetilde {u} (y) \| \leqslant \| u \| \| y \|\tag{5}
\]

对 E/N 的每个元素 \(y\) 成立。令 \(v = i \circ \widetilde{u}\) 。我们有 \(u = v \circ p\) ，并且

\[
((u)) = \inf_{\substack{y\in \mathrm{E / N}\\ y\neq 0}}\frac{\|v(y)\|}{\|y\|} = \inf_{\substack{y\in \mathrm{E / N}\\ \|y\| = 1}}\| v(y)\| .\tag{6}
\]

当 \(u\) 是零映射时，空间 E/N 缩减为 0，且有 \(((u)) = +\infty\) 与 \(\| u \| = 0\) 。当 \(u \neq 0\) 时，由 (3) 与 (4) 推出不等式

\[
0 \leqslant ((u)) \leqslant \| u \| <   + \infty .\tag{7}
\]

当 u 单射时，我们有

\[
((u)) = \inf_{\substack{x\in \mathrm{E}\\ x\neq 0}}\frac{\|u(x)\|}{\|x\|},\tag{8}
\]

且对一切 \(x \in \mathrm{E}\) ，我们有

\[
((u)) \| x \| \leqslant \| u (x) \| \leqslant \| u \| \| x \|.\tag{9}
\]

用 \(j\) 表示赋范空间 F 到其完备化 \(\widehat{\mathbf{F}}\) 中的典范单射。我们有 \(((u)) = ((j \circ u))\) 。

注记. --- 按定义，\(((u)) > 0\) 等价于双射线性映射 \(\widetilde{u}\) 是双连续的，亦即等价于 u 是严格态射（TG, III, p. 16）。此时我们有

\[
((u)) = \| \widetilde {u} ^ {- 1} \| ^ {- 1}.\tag{10}
\]

引理 2. --- 设 c 是实数。若 c \textless{} ((u))，则对每个元素 \(z \in \operatorname{Im}(u)\) ，存在 E 的元素 x，使得 \(u(x) = z\) 且 \(c \|x\| \leqslant \|z\|\) 。反之，若对每个 \(z \in \operatorname{Im}(u)\) 存在 \(x \in E\) ，使得 \(u(x) = z\) 且 \(c \|x\| \leqslant \|z\|\) ，则 \(c \leqslant ((u))\) 。

这是公式 (2)、公式 (5) 以及 \(u\) 的余范数定义的推论。

命题 6. --- 设 E 和 F 是赋范空间，\(u \in \mathcal{L}(\mathrm{E};\mathrm{F})\) 。用 B 表示 E 中范数 \(< 1\) 的元素所成的集。令

\[
\mathrm{P} = u (\mathrm{E}) - u (\mathrm{B}), \quad \mathrm{Q} = u (\mathrm{E}) - (\overline {{u (\mathrm{B})}} \cap u (\mathrm{E})).
\]

u 的余范数等于在 F 中从 0 到 P 的距离。若赋范空间 E 完备或 u 是严格态射，则 u 的余范数等于在 F 中从 0 到 Q 的距离。

设 N 是 u 的核；设 \(p: E \to E/N\) 是典范满射，并设 \(v: E/N \to F\) 是 u 在商空间上诱导的映射。集 \(p(B)\) 是 E/N 中范数 \textless{} 1 的元素所成的集。我们有

\[
\mathrm{P} = u (\mathrm{E}) - u (\mathrm{B}) = v (\mathrm{E} / \mathrm{N}) - v (p (\mathrm{B})) = v ((\mathrm{E} / \mathrm{N}) - p (\mathrm{B}))
\]

因为映射 \(v\) 是单射。换言之，P 由 F 中形如 \(v(y)\) （其中 \(y \in \mathrm{E/N}\) 且 \(\|y\| \geqslant 1\) ）的元素组成。用 \(d_{\mathrm{P}}\) 表示在 F 中从 0 到 P 的距离。我们有

\[
d_{\mathrm{P}} = \inf_{\substack{y\in \mathrm{E} / \mathrm{N}\\ \| y\| \geqslant 1}}\| v(y)\| = \inf_{\substack{y\in \mathrm{E} / \mathrm{N}\\ \| y\| = 1}}\| v(y)\| = ((u))
\]

由 (6) 即得。

设 \(u\) 是严格态射。设 \(\varepsilon > 0\) 。集 \(\varepsilon u(\mathrm{B})\) 是 \(u(\mathrm{E})\) 中 0 的一个邻域。\(u(\mathrm{B})\) 在 \(u(\mathrm{E})\) 中的闭包等于 \(\overline{u(\mathrm{B})} \cap u(\mathrm{E})\) 。它包含于集 \(u(\mathrm{B}) + \varepsilon u(\mathrm{B})\) ，后者等于 \((1 + \varepsilon)u(\mathrm{B})\) ，因为 \(u(\mathrm{B})\) 是凸的。于是有 \((1 + \varepsilon)\mathrm{P} \subset \mathrm{Q} \subset \mathrm{P}\) ，且在 F 中从 0 到 Q 的距离 \(d_{\mathrm{Q}}\) 满足不等式 \(d_{\mathrm{P}} \leqslant d_{\mathrm{Q}} \leqslant (1 + \varepsilon)d_{\mathrm{P}}\) 。由于这对每个 \(\varepsilon > 0\) 都成立，故有 \(d_{\mathrm{Q}} = d_{\mathrm{P}} = ((u))\) 。

设 \(u\) 不是严格态射，但赋范空间 E 完备。此时有 \(((u)) = 0\) 。\(u(\mathrm{B})\) 在 \(u(\mathrm{E})\) 中的闭包（它等于 \(\overline{u(\mathrm{B})} \cap u(\mathrm{E})\) ）不是 \(u(\mathrm{E})\) 中 0 的邻域（EVT, I, p. 17, 定理 1）。因此存在 Q 中任意接近 0 的点，从而 \(d_{\mathrm{Q}} = 0 = ((u))\) 。

命题 7. --- 设 E 是巴拿赫空间，F 是赋范空间。映射 \(u \mapsto ((u))\) （从 \(\mathcal{L}(\mathrm{E};\mathrm{F})\) 到 \(\overline{\mathbf{R}}\) 中）是上半连续的。

设 \(u \in \mathcal{L}(\mathrm{E};\mathrm{F})\) 。只需证明：对每个实数 \(c > ((u))\) ，由元素 \(v \in \mathcal{L}(\mathrm{E};\mathrm{F})\) （满足 \(((v)) < c\) ）所成的集是 \(u\) 的一个邻域。用 B 表示 E 中范数 \(< 1\) 的元素所成的集。由命题 6，存在 \(y \in \mathrm{E}\) ，使得 \(u(y) \notin \overline{u(\mathrm{B})}\) 且 \(\| u(y)\| < c\) 。距离 \(d\) （从 \(u(y)\) 到闭集 \(\overline{u(\mathrm{B})}\) ）是严格正的。由元素 \(v\) （\(\mathcal{L}(\mathrm{E};\mathrm{F})\) 中的、满足关系式 \(\| v(y)\| < c\) 与 \(\| u - v\|(1 + \| y\|) < d\) 者）所成的集 V 是 \(u\) 在 \(\mathcal{L}(\mathrm{E};\mathrm{F})\) 中的一个邻域。设 \(v \in \mathrm{V}\) 。对每个 \(x \in \mathrm{B}\) ，我们有

\[
\begin{array}{c} d \leqslant \| u (y) - u (x) \| \leqslant \| v (y) - v (x) \| + \| u - v \| (\| y \| + \| x \|) \\ <   \| v (y) - v (x) \| + d. \end{array}
\]

于是 \(v(y)\) 不属于 \(v(\mathrm{B})\) ，且由命题 6 有 \(((v)) \leqslant \| v(y) \| < c\) 。

命题 8. --- 设 E 是巴拿赫空间，F 是赋范空间。对每个 \(u \in \mathcal{L}(\mathrm{E};\mathrm{F})\) ，我们有 \(((u)) = ((^t u))\) 。

用 \(j\) 表示赋范空间 F 到其完备化 \(\widehat{\mathrm{F}}\) 中的典范单射。我们有 \(((u)) = ((j \circ u))\) 。由于线性映射 \(^t j: (\widehat{\mathrm{F}})' \to \mathrm{F}'\) 是双射且等距，也有 \(((^t u)) = ((^t (j \circ u)))\) 。因此只需就赋范空间 F 完备的情形证明本命题，我们在证明的其余部分即作此假设。

若 \(u\) 为零，则有 \(((u)) = ((^t u)) = +\infty\) 。若 \(u\) 不是严格态射，则 \(^t u\) 也不是（EVT, IV, p. 29, 推论 3），且有 \(((u)) = ((^t u)) = 0\) 。

现在设 u 是非零的严格态射。用 N 表示 u 的核，I 表示它的像，并考察 u 的典范分解：

\[
\mathrm{E} \xrightarrow {p} \mathrm{E/N} \xrightarrow {v} \mathrm{I} \xrightarrow {i} \mathrm{F}.
\]

\(^{t}u\) 的核是 I 在 F' 中的正交 I°（EVT, IV, p. 27, 命题 2），且 \(^{t}i\) 通过过渡到商空间诱导一个从 F'/I° 到 I' 上的等距 \(\iota\) （EVT, IV, p. 9, 命题 10）。此外，\(^{t}p\) 给出从 (E/N)' 到 N 在 E' 中的正交 N° 上的一个等距（EVT, IV, p. 8, 命题 9）。因此 \(^{t}u\) 的典范分解是

\[
\mathrm{F} ^ {\prime} \longrightarrow \mathrm{F} ^ {\prime} / \mathrm{I} ^ {\circ} \stackrel {v ^ {\prime}} {\longrightarrow} \mathrm{N} ^ {\circ} \longrightarrow \mathrm{E} ^ {\prime},
\]

其中 \(v^{\prime} = \pi \circ{}^{t}v\circ \iota\)。此时有

\[
\| (v ^ {\prime}) ^ {- 1} \| = \| (^ {t} v) ^ {- 1} \| = \| ^ {t} (v ^ {- 1}) \| = \| v ^ {- 1} \|
\]

（EVT, IV, p. 7, 命题 8），从而由公式 (10) 得 \(((u)) = ((^t u))\) 。

